\documentclass[11pt,a4paper]{article}
\usepackage[T1]{fontenc}
\usepackage{lmodern,amsmath,amssymb,amsthm,mathtools,booktabs,microtype}
\usepackage[margin=27mm]{geometry}
\usepackage[colorlinks=true,linkcolor=blue,urlcolor=blue,citecolor=blue]{hyperref}
\newtheorem{theorem}{Theorem}\newtheorem{lemma}[theorem]{Lemma}
\newtheorem{proposition}[theorem]{Proposition}
\newcommand{\Q}{\mathbb Q}\newcommand{\Z}{\mathbb Z}
\DeclareMathOperator{\Tr}{Tr}
\setlength{\emergencystretch}{2em}
\title{A defining polynomial in twenty-seven variables:\\
an existence theorem with a uniform degree bound}
\author{}\date{9 November 2026}
\begin{document}\maketitle
\begin{abstract}
Under the four arithmetic inputs of the specified elliptic source,
we prove that the integers are the rational input projection of an
integer polynomial in twenty-seven total variables and of degree
less than $2^{151}$. We give a computable family $F_L$ of such
polynomial specifications, indexed by a fixed odd integer coefficient
parameter $L$. Every sufficiently large $L$ has the required input
projection. The proof uses von Buhren's uniform Mordell--Lang
height theorem; its ineffective additive constant prevents us from
certifying a numerical sufficient value of $L$. Thus the variable
and degree bounds are proved, while no particular member is asserted
here to be a fully explicit defining polynomial. Exact source,
finite-algebra, and degree certificates accompany the family.
\end{abstract}

\section{Statement and arithmetic inputs}

All variable counts include the distinguished input. A polynomial
$F\in\Z[X,Y_1,\ldots,Y_m]$ defines $\Z$ over $\Q$ if
\[
 x\in\Z\quad\Longleftrightarrow\quad
 \exists\mathbf y\in\Q^m\quad F(x,\mathbf y)=0
 \qquad(x\in\Q).
\]
The integer $L$ below is fixed when a polynomial is constructed;
it is neither its input nor an existential variable.

\begin{theorem}\label{thm:main}
Assume the elliptic group and height assertion, the exponent
nineteen contact-height bound, the pole obstruction and
simultaneous-realization input, and the simultaneous prime-pattern input
specified in Section~\ref{sec:inputs}. There is a computable
family
\[
 F_L\in\Z[X,Y_1,\ldots,Y_{26}],\qquad
 L\ge3\text{ odd},
\]
of finite polynomial specifications with the following properties.
\begin{enumerate}
\item Every member has degree less than $2^{151}$.
\item There is a fixed real number $L_0$ such that every odd
      $L\ge L_0$ defines $\Z$ over $\Q$.
\item The proof does not determine an effective sufficient
      value of $L_0$, and does not give an effective
      coefficient-height bound for a member known to define $\Z$.
\end{enumerate}
\end{theorem}

The finite specification uses integer coefficient arithmetic,
determinants, resultants, and finite
products, all with explicit finite index ranges. It uniquely
specifies an ordinary polynomial, without treating the matrix
coordinates or algebra generators as quantified variables.

\subsection{The four inherited inputs}\label{sec:inputs}

Set
\[
 K=\Q(\sqrt2),\quad d=3-\sqrt2,\quad D=d^2=11-6\sqrt2,
 \quad E:y^2=x^3-Dx,
\]
\[
 P=\left(\frac{57}{8}-\frac72\sqrt2,
          -\frac{165}{8}+\frac{417}{32}\sqrt2\right).
\]
The construction uses the following assertions, with exactly the
scope established in the arithmetic dependency paper~\cite{Dependencies}.
\begin{enumerate}
\item $2E(K)=\Z P$, and $h(x([n]P))\ge n^2-6$ for every
      nonzero integer $n$.
\item For the five normalized contacts with
      $\infty,0,1,2401,-343/32$, the odd contact radical
      $M_5(q)$ outside primes at most $76176$ satisfies
      $h(q)\le H M_5(q)^{19}$ when $q$ is not a contact value,
      where
\[
 H=2^{110}210^{1800}\left(\prod_{p\le76176}p\right)^{19}+128.
\]
\item With $B=17b$, the pole system with preparation parameter
      zero is
\[
 p^2=B(1-e)(1+3e),\qquad
 Bq^2=(s^2-Be(B-b^2))(s^2+3Be(B-b^2)).
\]
      Its solvability excludes negative odd $v_\ell(b)$ for
      $\ell\notin\{2,3,17\}$. It is soluble whenever
      $0<b<17$, its reduced denominator is square, it is a
      square in $\Q_2,\Q_3,\Q_{17}$, and its valuations there
      are at least $3,1,2$, respectively. We retain the same
      restricted-converse and descent hypotheses supporting this
      assertion, and their simultaneous-realization construction.
      Lemma~\ref{lem:pole-chart} proves the nonzero-witness
      consequence of these inputs.
\item The simultaneous prime-pattern input permits the source's
      prescribed prime values, specified Gaussian generators, and reduction-order
      conditions, arbitrary fixed divisibility, and arbitrary
      fixed odd avoidance. It applies to the finite collection
      of comparison scales after the root points have been fixed.
\end{enumerate}
The prime-pattern input includes the stated application of Sun's
theorem. No new conjecture is introduced by the height theorem
used below: von Buhren's result~\cite{VonBuhren} is a published
theorem. Its effect here is ineffectivity of one fixed coefficient.

\section{The sum of four traces controls the four labels}

For nonzero integers $n_1,\ldots,n_4$ put
\[
 T(\mathbf n)=\Tr_{K/\Q}\left(\sum_{i=1}^4x([n_i]P)\right).
\]
The aggregate trace-height theorem~\cite{Aggregate} proves that
there are fixed constants $C>0$ and $J_0\ge0$ such that
\begin{equation}\label{eq:aggregate}
       \max_i|n_i|^2\le C h(T(\mathbf n))+J_0.
\end{equation}
The multiplicative constant comes from fixed geometric data;
the additive constant is ineffective.

For clarity, the essential hypotheses of this application are as
follows. After rescaling the two conjugate curves to
$E_0:y^2=x^3-x$, their source points are independent over
$\Z[i]$. The sum of their abscissas has no constant restriction
to an abelian translate through a tuple of nonzero multiples.
Within one block, coincident poles have tangent multipliers in
the rational ratios of their integer labels, so the leading
coefficient is a nonzero sum of positive rational squares.
Poles in different blocks cannot coincide by independence.
The closed graph of the trace function therefore avoids the
exceptional abelian translates at every source tuple. Applying
von Buhren's theorem to this fixed graph and its fixed finite-rank
subgroup gives~\eqref{eq:aggregate}, including a uniform additive
constant. The companion paper~\cite{Aggregate} supplies the full
proof over the algebraic closure, so translates defined only over
an extension field are included.

\section{The exact rational source}

We use the point charts, local tests, and denominator guards of
the established source~\cite{Baseline}. We describe their data
and the changes defining the present family.

There are nine points over $K$, named
\[
 A,S,T_1,Z,H_0,U_A,U_S,U_T,U_Z.
\]
For each point retain two rational coordinates $h_0,h_1$ and set
\begin{equation}\label{eq:directpoint}
 h=h_0+h_1\sqrt2,\qquad x=h^2,\qquad
 j^2=h^4-D,\qquad y=hj.
\end{equation}
This chart describes precisely the nonzero points of $\Z P$.
Indeed, let $R_0=(1+\sqrt2,2\sqrt2)$, so $2R_0=P$.
The group input implies that every $Q\in E(K)$ is $mR_0+T$
with $T\in E[2](K)$. For the standard two-isogeny descent map,
$\delta(R_0)=1+\sqrt2$ has opposite signs at the real embeddings,
whereas the torsion classes are $1,-D,d,-d$, with equal signs.
A square abscissa therefore forces $m$ even. Of these torsion
classes only $1$ is square: $-D,-d$ are negative and
$N_{K/\Q}(d)=7$ is not a rational square. Thus $Q\in\Z P$.
Conversely, $\delta(2R)=1$ gives a square abscissa for every
nonzero multiple of $P$. Neither $h=0$ nor $j=0$ is possible
in~\eqref{eq:directpoint}. A direct proof of the descent map
and all these assertions is given in~\cite{SquareCharts}.

The finite root $j$ is represented by one monic quartic,
without a finite choice of charts. Write
\[
 h^4-D=A+B\sqrt2,
 \quad A=h_0^4+12h_0^2h_1^2+4h_1^4-11,
 \quad B=4h_0h_1(h_0^2+2h_1^2)+6.
\]
Set
\[
 c=2A+B,\qquad d_1=-A-B,\qquad j'=(c+d_1\sqrt2)j.
\]
Writing $(c+d_1\sqrt2)^2(A+B\sqrt2)=A'+B'\sqrt2$ gives
\[
 A'=6A^3-9AB^2-4B^3,\qquad
 B'=-4A^3+6AB^2+3B^3.
\]
For $j=u+v\sqrt2$, direct expansion gives
\[
 j'=2(u^3-2v^3)+(4v^3-u^3)\sqrt2.
\]
Its rational component is nonzero, since $2$ is not a rational
cube and $j\ne0$. The multiplier has norm $2A^2-B^2$,
which is nonzero for rational $(A,B)\ne(0,0)$ because $2$
is not a rational square. Thus the reconstruction uses only
these two elementary rational obstructions. The general quadratic
extension version is proved in~\cite{QuadraticCharts}.
We may therefore use the single chart
\begin{equation}\label{eq:pointquartic}
 Y^4-4A'Y^2+8(B')^2=0,\quad Y\ne0,\qquad
 j'=Y/2+(B'/Y)\sqrt2,\qquad
 j=\frac{j'}{c+d_1\sqrt2}.
\end{equation}
The point-chart factors $Y$, $A'^2-2B'^2$, $2A^2-B^2$,
and the quartic derivative are automatically nonzero at every
rational retained-base specialization and every algebraic root.
Indeed, $A+B\sqrt2=h^4-D$ cannot vanish: taking norms would
imply $N(h)^4=49$, which has no rational solution. The homogeneous cubic
$-4A^3+6AB^2+3B^3$ has no nontrivial rational zero, because
$4T^3-6T-3$ is Eisenstein at $3$. Finally,
\[
 A'^2-2B'^2=(2A^2-B^2)^2(A^2-2B^2)
\]
and the discriminant of the quartic in~\eqref{eq:pointquartic} is
\[
                 2^{15}B'^2(2A^2-B^2)^4(A^2-2B^2)^2\ne0.
\]
The exact square-test proof is given in~\cite{SquarePredicate}.
These four factors per point may therefore be omitted from the
finite-algebra localization product. Every other denominator of the
affine point operations remains guarded.

The four root labels $(n,a,t,b)$ are nonzero integers; no parity
condition is needed for soundness. The simultaneous completeness
construction retains its original odd labels. The scalar shift
is now computed directly as $S-P$, in the affine difference
chart. Its excluded equal-abscissa cases are $S=\pm P$, both
avoided by that construction.

Put $Q_0=P$, $(Q_1,Q_2,Q_3,Q_4)=(A,S,T_1,Z)$,
$U_0=H_0$, and $V_i=U_i-Q_i$. Write $z=-x/y$ and $f=1/x$.
The truncated formal logarithms are
\[
 \ell_r(Q)=z(Q)\sum_{j=0}^{r/4-1}c_j(Df(Q)^2)^j,
 \quad c_0=1,\quad c_1=-2/5,\quad
 c_j=c_{j-1}\frac{4j-5}{4j+1}.
\]
Use precisions $(12,40,40,40)$, ratio orders $(10,39,39,39)$,
and scalar order $20$. With
$\beta=216X^4+450459$, the sixteen tests are the ten
formal-group tests $f(U_i),\sigma f(V_i)$, the four differences
between the paired logarithmic ratios, the direct target test
for $\beta$, and the direct Lucas polynomial
\[
 v^2-u(u-2)vw+(u-2)^2w^2-1
\]
in the three corresponding ratios. Their assigned orders are
$1$ for each of the first ten tests and
$10,39,39,39,10,39$ for the remaining six. These tests and
their construction are unchanged.

\subsection{Two comparison scales}

We use precisely the two positive rational
square scales
\begin{equation}\label{eq:scales}
 A_0=1,\qquad
 A_1=\left(\Tr_{K/\Q}\sum_{i=1}^4x(Q_i)\right)^2.
\end{equation}
They are nonzero: the root points are nonzero doubles and
their abscissas are positive in both real embeddings.
For the retained slope $\lambda$, each scale is inserted into
the same normalized comparison detector of~\cite{Baseline}:
\[
 \begin{gathered}
 t_r=A_r\lambda,\quad \delta_r=t_r-1,\quad
 F_r=\prod_{\theta\in\Theta}(t_r-\theta),\quad
 R_r=\frac{F_r}{A_r^5\delta_r^6},\\
 \Theta=\{0,-1/2,-1,1/2399,-32/407\}.
 \end{gathered}
\]
\[
 W_r=\prod_{j=1}^{16}\left(
 1+R_r\left(\frac{D_j}{\delta_r^{m_j+1}R_r^{m_j}}\right)^2\right),
 \qquad b_r=\frac{R_r}{\Tr(W_r^{-1})^2}.
\]
All displayed denominators and inherited guards are retained,
except for the automatic point-chart factors proved above.
Each $b_r$ receives a pole certificate.

For each pole put $B=17b$. We use the single rational chart
\[
 e=\frac{\tau^2-B}{\tau^2+3B},\qquad
 p=\frac{-4B\tau}{\tau^2+3B},\qquad
 \tau^2+3B\ne0,
\]
with retained parameters $\tau,s$, and require the remaining quadratic
pole root $q$ to be nonzero. The following lemma proves completeness
of this restriction under the inherited arithmetic inputs.

\begin{lemma}\label{lem:pole-chart}
Suppose that $b$ has square reduced denominator,
\[
 0<b<17,\qquad b\in\Q_2^{\times2}\cap\Q_3^{\times2}
 \cap\Q_{17}^{\times2},\qquad
 v_2(b)\ge3,\quad v_3(b)\ge1,\quad v_{17}(b)\ge2.
\]
The inherited pole-realization proof supplies a certificate with
$pqs\ne0$ and $e\ne1$. Consequently the single chart above,
with $q\ne0$, suffices. The restricted predicate has the same
pole obstruction, and adds no arithmetic hypothesis.
\end{lemma}
\begin{proof}
Set $V=B(B-b^2)$. Its positivity and its local square conditions
are precisely those verified in the zero-preparation argument
of~\cite{Dependencies}. Let $\beta_0,\theta$ be the positive
squarefree representatives of $B,V$. Their supports are disjoint,
$\beta_0\equiv17\pmod{24}$, $\theta\equiv1\pmod{24}$, and
both mutual local-square conditions hold.

The inherited simultaneous-realization proof in~\cite{Glazer},
with source identifier
\begin{center}\small
\texttt{lem:original-seven-simultaneous-realization},
\end{center}
constructs points on $E_n:y^2=x(x-n)(x+3n)$ and $E_{n\theta}$ with nonzero
abscissas of squareclass $\beta_0$. Both are outside two-torsion.
Indeed, the construction includes a fresh prime of $n$ outside
$\beta_0\theta$, so $\beta_0$ differs in squareclass from $n$
and $n\theta$; its positive squareclass also differs from their
negative multiples by $-3$. These are the only nonzero two-torsion
abscissas. This non-torsion assertion is also explicit in the
source descent proof. Thus both ordinates are nonzero.

Writing $V=\theta c^2$ and rescaling as in that proof gives
\[
       (B,y_1)\in E_{e_0}(\Q),\qquad
       (Bs^2,y_2)\in E_{e_0V}(\Q),\qquad sy_1y_2\ne0.
\]
The substitutions $e=e_0/B$, $p=y_1/B$, $q=y_2/(B^2s)$
give exactly the two pole equations and $pqs\ne0$. The first
conic excludes $e=1$. Set $\tau=p/(e-1)$. Its equation gives
\[
       \tau^2+3B=\frac{4B}{1-e}\ne0,
\]
and then gives the displayed reconstruction of $e,p$.
Every restricted witness is an original pole witness, so its
obstruction is unchanged. The construction establishes completeness
for every input satisfying the displayed conditions. No equivalence
of the two pole predicates for arbitrary rational $b$ is asserted.
\end{proof}

\subsection{A linear scalar normalization}

The fixed scalar normalization is
\begin{equation}\label{eq:normalizer}
 \begin{gathered}
 \ell=76249,\quad
 r=\beta(\beta-1)(\beta-2401)(32\beta+343)(2\beta-1),\\
 R_L=\frac{102^2\ell^L r}{(1+C_Lr)^2},\quad
 C_L=17(\ell^{(L+1)/2}+1).
 \end{gathered}
\end{equation}
Here $L\ge3$ is odd. The scalar norm tree and its six rational
component tests are unchanged: they concern
$f(H_0),f(S-P),z(S-P)/z(H_0)$ with orders
$1,1,1,1,20,20$. Its output $W$ supplies the scalar pole
input $R_LW^2$.

The finite residue certificate verifies that the five contact
values are not attained by $216x^4+450459$ modulo $\ell$,
and that $1+17r$ never vanishes there. Consequently, for every
rational input,
\begin{equation}\label{eq:valuation}
 v_\ell(R_L)=
 \begin{cases}
 L,&v_\ell(X)\ge0,\\
 L-20v_\ell(X),&v_\ell(X)<0.
 \end{cases}
\end{equation}
These valuations are positive and odd. At a consumed contact
prime different from $\ell$, the denominator is a unit and the
contact valuation is unchanged. For integral $X$, the reduced
denominator of $R_L$ is square, it has the required local
squareclasses and depths at $2,3,17$, and
\[
 0<R_L<\frac{36}{\ell r}<1.
\]
The scalar norm argument therefore gives
\begin{equation}\label{eq:depth}
 v_\ell(a-1)\ge19L+1.
\end{equation}
The normalization calculation in~\cite{Aggregate}, together with
the finite-field avoidance certificate~\cite{Normalizer}, proves
these statements and the preparation-zero realization. The supplied
checker independently verifies every residue modulo $76249$
and the normalization identity as an exact polynomial identity.

\section{Soundness}

The affine guard for $S-P$ excludes $a=1$. Let $c=|a-1|>0$ and
$B=\max(|n|,|a|,|t|,|b|)$. Equation~\eqref{eq:depth} and $L\ge3$
give $c\ge\ell^{19L+1}$ and therefore $B\ge c-1>4$.
In particular, the elementary residual bound $|\Delta|\le5B^4$
used below is valid. If the two contact radicals at the
scales~\eqref{eq:scales} have their $k$th powers dividing
$N\ne0$, the contact-height theorem gives
\[
 h(\lambda),\ h(A_1\lambda)
       \le(H+\log3)|N|^{19/k}.
\]
Since $h(A_1)=2h(\Tr\sum_i x(Q_i))$, the height inequality
for a quotient and~\eqref{eq:aggregate} give constants
$K\ge1$, $J\ge0$, independent of $X,L$, such that
\begin{equation}\label{eq:commonheight}
                 B^2\le K|N|^{19/k}+J.
\end{equation}
The local formal-logarithm arguments of the source are
unchanged; hence, with
$\Delta=t^2-a(a-2)tb+(a-2)^2b^2-1$, they imply
\begin{align}
 h(\beta)&\le Hc^{19/20},\label{eq:betaheight}\\
 \Delta\ne0&\Longrightarrow
 B^2\le K(5B^4)^{19/39}+J,\label{eq:delta}\\
 \beta\ne n&\Longrightarrow
 B^2\le K(2e^{h(\beta)}B)^{19/10}+J.\label{eq:beta}
\end{align}

Choose a fixed threshold $T$ greater than each of
\begin{equation}\label{eq:threshold}
 \begin{gathered}
 16,\qquad 2\sqrt{2J}+2,\qquad
 2(2K5^{19/39})^{39/2}+2,\\
 (128H)^{20},\qquad
 8\bigl(10\log(2K)+19\log2\bigr).
 \end{gathered}
\end{equation}
For $c>T$, one has $B\ge c-1>c/2$ and $B^2>2J$.
If $\Delta\ne0$, equation~\eqref{eq:delta} gives
$B^{2/39}\le2K5^{19/39}$, contrary to~\eqref{eq:threshold}.
Thus $\Delta=0$. The elementary Lucas growth lemma proved
in~\cite{Baseline} then gives $\log B>c-4$.
If $\beta\ne n$, equation~\eqref{eq:beta} gives
\[
 \log B\le10\log(2K)+19\log2+19h(\beta)
       <\frac c8+\frac{19c}{128}<c-4,
\]
using~\eqref{eq:betaheight} and the threshold. This contradiction
proves $\beta=n\in\Z$. Finally the elementary input-map
lemma of~\cite{Baseline} says that
$216X^4+450459\in\Z$ for rational $X$ if and only if
$X\in\Z$.

By~\eqref{eq:depth}, every source solution has
$c\ge\ell^{19L+1}$. Therefore every odd $L$ with
$\ell^{19L+1}>T$ satisfies the soundness conclusion. The
constant $T$ is fixed but ineffective because $J$ is ineffective.

\section{Completeness}

Fix an integer input $X$ and a fixed odd $L$. All numerator
and denominator primes of~\eqref{eq:normalizer} form a finite
set. The source realization first chooses the odd root labels
with the scalar label $a-1$ sufficiently divisible and deep at
these primes, and with the corresponding points in the required
real neighborhoods. The explicit Lucas construction then gives
odd nonzero $t,b$ with $\Delta=0$ and arbitrarily large
growth. The input label is chosen to be $n=\beta$.

At this stage the two scales in~\eqref{eq:scales} are fixed
positive rational squares. Apply the same simultaneous
prime-pattern theorem to these two scales. It chooses the slope
and contact primes with the prescribed conjugate reduction
orders while avoiding the previously fixed odd modulus. The
subsequent Chinese remainder choice of the auxiliary index $h$
and the points $U_i=[hn_i]P$ supplies the sixteen local tests
to their assigned precisions and preserves the fixed scalar
requirements and real neighborhoods. This is precisely the
simultaneous construction of~\cite{Dependencies,Baseline}; its
prime-pattern input accepts any fixed finite collection of the
positive square scales and any fixed finite avoidance modulus.

Each of the two comparison inputs is thereby prepared for the
zero-preparation pole theorem. The scalar input is prepared as
well: $R_L$ has square reduced denominator and the required
local squareclasses, the six normalized scalar tests can be made
arbitrarily small at the fixed finite places, and the scalar
norm output can be made arbitrarily close to $1$ in the real
absolute value. Thus $0<R_LW^2<17$, with the requisite
denominator and local conditions. The three pole certificates
exist by the inherited completeness input. The direct point chart~\eqref{eq:pointquartic} covers all nine
chosen points; its multiplier needs no further choice.
Lemma~\ref{lem:pole-chart} applies independently to each pole input.
Its parameters belong only to that pole, so the certificates can
be chosen simultaneously without any common-parameter assumption.
All pole roots are nonzero and the single conic chart applies.
All remaining inherited guards are maintained by the same realization.

This proves completeness for every fixed odd $L\ge3$.
Together with the preceding section it proves the arithmetic
assertion for every sufficiently large such $L$.

\section{Elimination of finite roots and the uniform degree}

The retained coordinates and finite roots are as follows.
\begin{center}
\begin{tabular}{lr}
\toprule
Retained rational coordinates & Count\\
\midrule
Input & 1\\
Nine square-root-abscissa pairs & 18\\
Comparison slope & 1\\
Two comparison-pole parameter pairs & 4\\
Scalar-pole parameter pair & 2\\
\midrule
Total before finite-algebra elimination & 26\\
\bottomrule
\end{tabular}
\end{center}
There are nine quartic point roots and three quadratic pole roots.
The finite-algebra rank is $4^9 2^3=2^{21}$. Each pole root is
already guarded nonzero in the source, so every zero stratum is
inconsistent and is discarded. The first conics use one chart.
The direct-scaled finite-algebra theorem of~\cite{Weighted} replaces
all twelve roots by one rational coordinate and retains every
nonautomatic nonzero denominator condition. This gives twenty-seven
total polynomial variables.

Let $n_0$ be the rank of the finite algebra and let $m$ be its
number of monic generators $z_1,\ldots,z_m$. Let $G$ be the product
of the remaining nonzero guards. For each integer
\[
 0\le s\le m n_0(n_0+1)/2,\qquad
 \lambda_s=1+\sum_{i=1}^{m}s^i z_i,
\]
form the multiplication matrix and its characteristic polynomial
\[
 H(V)=\det(I-VM_{G\lambda_s})=\sum_{k=0}^{n_0}a_kV^k.
\]
There are no closed residual equations in this source. For every
$1\le d\le n_0$, define
\[
 H_d(V)=\sum_{k=0}^d a_kV^k,\qquad
 g_d=\operatorname{Res}_V(H_d,H_d'),
\]
using Sylvester degrees $d,d-1$. The factor for this case is
\[
 Q_{s,d}(U)=
 \begin{cases}
 H_d(g_dU)^2+\displaystyle\sum_{k=d+1}^{n_0}a_k^2,&d<n_0,\\[1mm]
 H_{n_0}(g_{n_0}U),&d=n_0.
 \end{cases}
\]
The defining polynomial is the product of these factors over all
specified $s,d$. Every operation is integer coefficient arithmetic
with explicit finite index ranges.

The direct-scaled finite-algebra theorem of~\cite{Weighted} proves
that this product has exactly the rational base projection of the
guarded source. Indeed, $H_d(0)=1$ excludes $g_d=0$ at every zero
of a factor, and the resultant then makes the full characteristic
polynomial squarefree of degree $d$. Its rational root gives a
one-dimensional rational eigenspace for all commuting multiplication
matrices, hence a guarded rational point. Conversely one of the
specified separating elements makes all active eigenvalues distinct
and nonzero. Inactive nonreduced fibers, complex residue factors,
and drops in degree are included in that proof.

For the quantitative bounds, give each retained variable weight one
and each monic generator its certified weight. If the multiplication
element $G\lambda_s$ has weight at most $w$, then
$\deg a_k\le kw$: the basis-weight differences cancel in every
principal minor. The resultant has coefficient index weight $d^2$.
Consequently, with
\[
 S=1+\frac{m n_0(n_0+1)}2,\qquad
 C(n_0)=2\sum_{d=1}^{n_0-1}\max(d^3+d,n_0)+(n_0^3+n_0),
\]
we have
\[
                      \deg F\le S\bigl(wC(n_0)+n_0^2\bigr).
\]
The complete degree and coefficient-norm proofs are supplied
in~\cite{Weighted}.

Here $m=12$, $n_0=2^{21}$, and the generator weights are
\[
 \underbrace{6,\ldots,6}_{9\text{ entries}},
 1568138,\ 1935231,\ 116318.
\]
The guard has weight $4423559$, giving $w=6358790$.
There are exactly $55340258609409818624$ factors.
The independent integer bound is
\[
 D_*=1622838857795564438964780120527413199713157888
       <2^{151}.
\]
It is uniform in $L$. The source uniform-degree checker replaces
$\ell^L$ and $\ell^{(L+1)/2}$ by independent coefficient symbols,
each of degree zero in the retained variables. Integer specialization
cannot increase any source weight, and the same weighted
principal-minor calculation applies for every fixed odd $L$.
No coefficient-height bound for a member certified to exceed the
ineffective sufficient threshold is inferred.

\section{Reproduction and the status of a numerical example}

The directory \texttt{trace27\_nontorsion} contains the exact guarded
source, its builder, independent source checks, and its uniform-degree
certificate. The source command
\begin{verbatim}
python3 reproduce.py --anchor-exponent 137731
\end{verbatim}
reconstructs and checks the numerical source. The exponent is a fixed
coefficient choice and does not add a quantified variable. Its source
has $7369$ coefficient nodes; integer coefficient conversion gives
$3032$ nodes. Source verification covers nine quartic point charts,
three rational-square clearings, all $104$ atomic source guard factors,
and the $36$ automatic factors omitted from the guard. The scalar
normalizer is checked at all $76249$ residue classes.

The directory \texttt{weighted\_elimination} contains the compiler and
independent checker:
\begin{verbatim}
build_weighted_specification.py
check_weighted_specification.py
\end{verbatim}
The final specification is in \texttt{weighted27}. The compiler uses the verified source
coefficient graph and writes exactly the determinant and finite product
defined above. The checker evaluates the integer bound sums literally
and verifies small exact algebras exercising singular and inactive
fibers, complex residue factors, selector collisions, and degree drops.
All programs require only the Python standard library.

For the numerical exponent $137731$, the resulting coefficient height
is bounded by $2^{C_*}$ with
\[
 C_*=3287373449762048678524372378187233080614463290580<2^{162}.
\]
This is a bound for a specified algebraic family member. No argument
here certifies that this numerical exponent exceeds the ineffective
sufficient threshold. The defining-polynomial assertion is the
existence theorem for every sufficiently large odd exponent, with
the uniform degree bound proved above.

\begin{thebibliography}{99}
\bibitem{Weighted}
\emph{Weighted bounds for eliminating finite algebraic choices},
companion paper, 9 November 2026,
\nolinkurl{weighted_finite_algebra_20261109_v12.tex}.

\bibitem{Aggregate}
\emph{Aggregate trace heights of elliptic multiples},
companion paper, 8 November 2026.
\bibitem{Baseline}
\emph{An explicit polynomial in 29 variables over the rationals},
the established 29-variable source and guarded finite-algebra
construction, version of 8 November 2026,
\nolinkurl{few_variables_20261108_v8.tex}.
\url{https://github.com/JoseBrox/diophantine-z-q_improvements}.
\bibitem{Dependencies}
\emph{Arithmetic dependencies of the construction with 29 variables},
version of 8 November 2026,
\nolinkurl{arithmetic_dependencies_20261108_v8.tex}.
\url{https://github.com/JoseBrox/diophantine-z-q_improvements}.
\bibitem{Normalizer}
\emph{Scalar contact normalization by finite-field avoidance},
companion scalar-normalization paper, 8 November 2026.
\bibitem{SquareCharts}
\emph{Square abscissas and direct charts for the elliptic subgroup},
companion paper, 8 November 2026.
\bibitem{QuadraticCharts}
\emph{A single quartic chart for square roots in a quadratic extension},
companion paper, 8 November 2026,
\nolinkurl{quadratic_square_chart_20261108_v9.tex}.
\bibitem{VonBuhren}
J.~von Buhren,
\emph{Borne de hauteur semi-effective pour le probl\`eme de
Mordell--Lang dans une vari\'et\'e ab\'elienne},
J.~Th\'eor. Nombres Bordeaux \textbf{29} (2017), no.~1,
307--320, Theorem~1.1.
\url{https://www.numdam.org/item/10.5802/jtnb.981.pdf}.
\bibitem{Glazer}
E.~Glazer, \emph{Diophantine definition of the integers over the
rationals}, Version~9, \nolinkurl{manuscripts/quartic_v9.tex},
commit \nolinkurl{14e8f601fe0d9173d4a5f084aedc3e07066a7688}.
\url{https://github.com/elliotglazer/diophantine-z-q}.
The inherited restricted-converse correction is identified
in~\cite{Dependencies}.
\bibitem{SquarePredicate}
\emph{Optimal polynomial square tests in quadratic number fields},
companion paper, 8 November 2026,
\nolinkurl{quadratic_square_predicates_20261108_v10.tex}.
\end{thebibliography}
\end{document}
